% Lösungen Körper (zusammenbau v0.4, Heft)
\einheitenkopf{Körper}

\erg{9}{a) Zylinder \quad b) Quader \quad c) $8$; $12$; $6$ (Ecken; Kanten; Flächen)}
\erg{10}{a) Prisma \quad b) Prisma \quad c) Pyramide – ein Quadrat als Grundfläche, vier Dreiecke laufen in einer Spitze zusammen \quad d) Diagonalen der Grundfläche zeichnen, ihren Schnittpunkt gestrichelt mit der Spitze verbinden: das ist die Höhe, $2{,}4$ m; eine Grundkante mit $3{,}6$ m beschriften – nicht eine Seitenkante und nicht die Höhe einer Seitenfläche \quad e) Diagonalen der Grundfläche zeichnen, ihren Schnittpunkt gestrichelt mit der Spitze verbinden: das ist die Höhe, $1{,}9$ m; eine Grundkante mit $2{,}8$ m beschriften – nicht eine Seitenkante und nicht die Höhe einer Seitenfläche \quad f) Diagonalen der Grundfläche zeichnen, ihren Schnittpunkt gestrichelt mit der Spitze verbinden: das ist die Höhe, $7$ cm; eine Grundkante mit $9$ cm beschriften – nicht eine Seitenkante und nicht die Höhe einer Seitenfläche}

10d)

\pyramide{3.5}{3.5}{3}{$3{,}6$ m}{}{$2{,}4$ m}


10e)

\pyramide{3.5}{3.5}{2.5}{$2{,}8$ m}{}{$1{,}9$ m}


10f)

\pyramide{3.5}{3.5}{3.2}{$9$ cm}{}{$7$ cm}

\erg{11}{a) M – in der Reihe D, N, M liegt genau eine Fläche dazwischen; N grenzt an D an \quad b) M – in der Reihe M, N, D liegt genau eine Fläche dazwischen; N grenzt an D an \quad c) $6 + 6 = 12$ Kanten ($6$ Grundkanten, $6$ Seitenkanten; nicht $7$ Ecken) \quad d) $3 + 3 = 6$ Kanten ($3$ Grundkanten, $3$ Seitenkanten; nicht $4$ Ecken)}
\erg{12}{a) Grundkreis auf dem Boden, er berührt alle vier Seitenwände (Durchmesser $50$ cm $=$ Kartonbreite), im Schrägbild als Ellipse; Spitze in der Mitte des Deckels (Schnitt der Diagonalen); zwei Mantellinien von der Spitze zum Kreisrand \quad b) Grundkreis auf dem Boden, er berührt alle vier Seitenwände (Durchmesser $36$ cm $=$ Kartonbreite), im Schrägbild als Ellipse; Spitze in der Mitte des Deckels (Schnitt der Diagonalen); zwei Mantellinien von der Spitze zum Kreisrand \quad c) Grundkreis auf dem Boden, er berührt alle vier Seitenwände (Durchmesser $18$ cm $=$ Kartonbreite), im Schrägbild als Ellipse; Spitze in der Mitte des Deckels (Schnitt der Diagonalen); zwei Mantellinien von der Spitze zum Kreisrand}
\erg{13}{a) Volumen \quad b) $6 \cdot 3 \cdot 2 = 36$ cm³ \quad c) $5 \cdot 5 \cdot 5 = 125$ cm³}
\erg{14}{a) $8 \cdot 8 \cdot 8 = 512$ cm³ (nicht $8 \cdot 8 = 64$, nicht die Oberfläche $384$ cm²) \quad b) $2 \cdot 2 \cdot 2 = 8$ dm³ (nicht $2 \cdot 2 = 4$, nicht die Oberfläche $24$ dm²)}
\erg{15}{a) Grundfläche: eines der beiden gleichen Dreiecke (die Deckfläche ist das andere); Höhe: eine Kante zwischen den Dreiecken, der Abstand von Grund- und Deckfläche – auch wenn das Prisma liegt \quad b) $18 \cdot 7 = 126$ cm³ \quad c) Grundfläche $7 \cdot 3 = 21$ cm²; $21 \cdot 5 = 105$ cm³}
\erg{16}{a) $2\,340 \cdot 85 = 198\,900$ cm³ (die gegebene Trapezfläche nutzen, nicht neu berechnen) \quad b) $418 \cdot 72 = 30\,096$ cm³ (die gegebene Trapezfläche nutzen, nicht neu berechnen) \quad c) Deckfläche $\tfrac{1}{2} \cdot 1{,}8 \cdot 1{,}56 = 1{,}404$ m²; Volumen $1{,}404 \cdot 0{,}75 \approx 1{,}05 \approx 1{,}1$ m³ \quad d) Deckfläche $\tfrac{1}{2} \cdot 0{,}7 \cdot 0{,}61 = 0{,}2135$ m²; Volumen $0{,}2135 \cdot 2{,}4 \approx 0{,}51 \approx 0{,}5$ m³}
\erg{17}{a) Durchmesser; Radius $4$ cm \quad b) $\pi \cdot 5^2 \cdot 8 \approx 628{,}3$ cm³ \quad c) $r = 4$ cm; $\pi \cdot 4^2 \cdot 5 \approx 251{,}3$ cm³}
\erg{18}{a) $\sqrt{600 : (\pi \cdot 9)} \approx 4{,}6$ cm (erst teilen, dann die Wurzel) \quad b) $\sqrt{2\,000 : (\pi \cdot 15)} \approx 6{,}5$ cm (erst teilen, dann die Wurzel) \quad c) $\sqrt{1\,500 : (\pi \cdot 18)} \approx 5{,}2$ cm (erst teilen, dann die Wurzel) \quad d) $\pi \cdot 3{,}5^2 \cdot 9 \approx 346{,}4$ cm³ $\approx 346$ ml – das sind etwa $350$ ml \quad e) $\pi \cdot 3{,}6^2 \cdot 8{,}5 \approx 346{,}1$ cm³ – nein, knapp unter $350$ ml \quad f) $\pi \cdot 5{,}2^2 \cdot 11{,}5 \approx 976{,}9$ cm³ $\approx 0{,}98$ l – etwa ein Liter \quad g) $2{,}5$ l $= 2\,500$ cm³; Höhe $= 2\,500 : (\pi \cdot 7^2) \approx 16{,}2$ cm \quad h) $0{,}75$ l $= 750$ cm³; Höhe $= 750 : (\pi \cdot 5^2) \approx 9{,}5$ cm}
\erg{19}{a) Quader (Wohnteil) und Dreiecksprisma (Dach) \quad b) $8 \cdot 4 \cdot 2 + 8 \cdot 2 \cdot 2 = 64 + 32 = 96$ dm³ \quad c) $(7 \cdot 4 + \tfrac{1}{2} \cdot \pi \cdot 2^2) \cdot 1{,}6 \approx 54{,}9$ m³ – die Tiefe gilt für Rechteck und Halbkreis, darum die Klammer}
\erg{20}{a) $37$ cm × $37$ cm × $56$ cm, weil der Durchmesser $= 2 \cdot 18 = 36$ cm in Länge und Breite passen muss und die Höhe $55$ cm \quad b) $26$ cm × $26$ cm × $11$ cm, weil der Durchmesser $= 2 \cdot 12{,}5 = 25$ cm in Länge und Breite passen muss und die Höhe $9$ cm \quad c) $41$ cm × $41$ cm × $71$ cm, weil der Durchmesser $= 2 \cdot 20 = 40$ cm in Länge und Breite passen muss und die Höhe $70$ cm \quad d) Zylinder $\pi \cdot 18{,}5^2 \cdot 36 \approx 38\,708$ cm³; Kegel $\tfrac{1}{3} \cdot \pi \cdot 18^2 \cdot 35 \approx 11\,875$ cm³; Rest $\approx 26\,832$ cm³ \quad e) Hülle $\pi \cdot 4{,}2^2 \cdot 11{,}5 \approx 637{,}3$ cm³; Kegel $\tfrac{1}{3} \cdot \pi \cdot 4^2 \cdot 11 \approx 184{,}3$ cm³; Rest $\approx 453{,}0$ cm³}
